\scanpage{89}
- \mspage{IV-1009} -
\uline{On a alors les implications suivantes} :
\diagram{\begin{tikzcd}[ampersand replacement=\&,column sep=small,row sep=small]
\text{a)}\arrow[r,Rightarrow]\arrow[dr,Rightarrow]
\&\text{e)}\arrow[r,Rightarrow]\arrow[dr,Rightarrow]
\&\text{b)}\arrow[r,Leftrightarrow]
\&\text{c)}\arrow[r,Leftrightarrow]
\&\text{h)}\&\\
\&\text{g)}\arrow[r,Rightarrow]
\&\text{d)}\arrow[r,Leftrightarrow]
\&\text{d')}\arrow[r,Leftrightarrow]
\&\text{f)}\arrow[r,Leftrightarrow]
\&\text{f')}
\end{tikzcd}}
\uline{En outre, lorsque tout module quotient d'un sous-module d'un}
$\operatorname{gr}_k(M)$
\uline{est séparé pour la topologie} $\underline J$-\uline{préadique}
(ce qui sera le cas lorsque
$A$ est noethérien, que les $f_i$ appartiennent au radical de $A$ et que
les $\operatorname{gr}_k(M)$ sont des $A$-modules de type fini),
\uline{alors les conditions
a) à h) sont toutes équivalentes}.

Notons que d'après (19.5.5), il résulte de l'hypothèse que $f$ est
$\operatorname{gr}_\bullet(N)$-régulière, que $\psi_N$ est \uline{bijectif};
comme $\unsure{\operatorname{gr}''(j)|\operatorname{gr}'(j)}$ est injectif,
l'équivalence de d) et f) résulte du diagramme (19.6.2.1), ainsi
que celle de d') et f'); par ailleurs f) et f') sont trivialement
équivalentes, ce qui montre l'équivalence de d),d'),f) et f').
Comme on sait déjà que $\psi_M$ est surjectif et $\psi_N$ bijectif, l'équivalence
de b) et c) découle aussi du diagramme (19.6.2.1) par un cas
particulier du lemme des cinq.

La relation $\psi'_R=\gamma\circ\psi_R$\inserted{ (19.6.2)}
et le fait que $\psi_R$ est surjectif montre\inserted{nt} que la
condition c) équivaut à dire que $\gamma$ est surjectif, ce qui équivaut
encore à la condition h); on a donc prouvé l'équivalence de b),c)
et h).

Il est trivial que e) équivaut à la \deleted{\unsure{conjonction}} conjonction de c)
et d), donc aussi de b) et d). Notons maintenant que a) implique
b) \deleted{\illegible{1}} et que b) implique a) sous l'hypothèse de séparation
(19.5.5). D'autre part, les implications a)$\Rightarrow$g)$\Rightarrow$\deleted{\unsure{f)}} f'),
et
l'implication f')$\Rightarrow$a) sous l'hypothèse de séparation, résultent de
(19.4.3). On a donc prouvé que a) entraîne e) (équivalente à la conjonction
de b) et f')), et aussi que toutes les conditions sont équivalentes
sous l'hypothèse de séparation. CQFD.

\uline{Corollaire} (19.6.4) .- \uline{Soient} $\underline K$ \uline{un idéal de} $A$,
$\underline L=\underline J+\underline K$. \uline{Supposons
que la filtration} $(N_k)$ (\uline{donc aussi} $(M_k)$)
\uline{soit la filtration} $\underline K$-\uline{\unsure{préadi-}}
\mspage{88}
